% ================================================================== Chapitre 5
\chapter{Deuxième principe : évolution et équilibre}
\label{ch:evolution}

On établit ici ce qui était admis en première année : le critère d'évolution d'un système chimique et la loi d'action des masses. Tout découle d'une seule idée : à $T$ et $P$ fixées, l'enthalpie libre d'un système ne peut que diminuer.

\noindent\emph{Cadre de tout le chapitre} : système fermé, siège d'une réaction unique, maintenu à température et pression fixées ($T = T_{\text{ext}}$, $P = P_{\text{ext}}$), où seules les forces de pression travaillent.

\section{Création d'entropie et enthalpie libre de réaction}

\begin{resultat}{Création d'entropie due à la réaction}
\[ \boxed{\delta S_c = -\frac{\Delta_rG\,\dd\xi}{T} \geq 0} \qquad\Longleftrightarrow\qquad \dd G = -T\,\delta S_c \leq 0 \]
\end{resultat}

\begin{demo}
Sur une évolution élémentaire :
\begin{itemize}
  \item premier principe (monobare, forces de pression seules) : $\dd H = \delta Q$ ;
  \item deuxième principe (monotherme, à la température $T$ du thermostat) : $\dd S = \dfrac{\delta Q}{T} + \delta S_c$ avec $\delta S_c \geq 0$ ;
  \item $T$ étant fixée : $\dd G = \dd H - T\,\dd S = \delta Q - \delta Q - T\,\delta S_c = -T\,\delta S_c$.
\end{itemize}
Par ailleurs, à $T$ et $P$ fixées, $\dd G = \sum_i\mu_i\,\dd n_i = \sum_i\mu_i\nu_i\,\dd\xi = \Delta_rG\,\dd\xi$. En égalant : $T\,\delta S_c = -\Delta_rG\,\dd\xi$.
\end{demo}

L'enthalpie libre joue donc, à $T$ et $P$ fixées, le rôle d'un \textbf{potentiel thermodynamique} : elle décroît au cours de toute évolution spontanée et elle est \emph{minimale à l'équilibre}. L'irréversibilité de la transformation est entièrement due à la réaction elle-même : $\delta S_c$ est non nulle tant que $\Delta_rG \neq 0$.

\section{Critère d'évolution spontanée}

\begin{resultat}{Sens d'évolution à $T$ et $P$ fixées}
Le deuxième principe impose $\Delta_rG\,\dd\xi \leq 0$ :
\begin{itemize}
  \item $\Delta_rG < 0 \;\Rightarrow\; \dd\xi > 0$ : évolution dans le \textbf{sens direct} ($\rightarrow$) ;
  \item $\Delta_rG > 0 \;\Rightarrow\; \dd\xi < 0$ : évolution dans le \textbf{sens indirect} ($\leftarrow$) ;
  \item $\Delta_rG = 0$ : \textbf{équilibre chimique}, aucune évolution.
\end{itemize}
\end{resultat}

\begin{pieges}
\begin{itemize}
  \item C'est le signe de $\Delta_rG$, \emph{pas} celui de $\Delta_rG^\circ$, qui donne le sens d'évolution. $\Delta_rG^\circ$ ne dépend que de $T$ ; $\Delta_rG$ dépend aussi de la composition actuelle.
  \item Le critère dit si une évolution est \emph{possible}, pas si elle est \emph{rapide} (exemple du diamant, chapitre~\ref{ch:potentiel}).
\end{itemize}
\end{pieges}

\section{Expression de l'enthalpie libre de réaction}

\begin{resultat}{Enthalpie libre de réaction et quotient de réaction}
\[ \boxed{\Delta_rG = \Delta_rG^\circ + RT\ln Q_r} \qquad\text{avec}\qquad Q_r = \prod_i a_i^{\,\nu_i} \]
\end{resultat}

\begin{demo}
On injecte $\mu_i = \mu_i^\circ + RT\ln a_i$ dans $\Delta_rG = \sum_i\nu_i\mu_i$ :
\[ \Delta_rG = \underbrace{\sum_i\nu_i\mu_i^\circ}_{\Delta_rG^\circ} + RT\sum_i\nu_i\ln a_i = \Delta_rG^\circ + RT\ln\Big(\prod_i a_i^{\,\nu_i}\Big) = \Delta_rG^\circ + RT\ln Q_r. \]
\end{demo}

\section{Constante d'équilibre et loi d'action des masses}

\begin{definition}[constante thermodynamique d'équilibre]
La \textbf{constante thermodynamique d'équilibre} d'une réaction à la température $T$ est
\[ \boxed{K^\circ(T) = \exp\left(-\frac{\Delta_rG^\circ(T)}{RT}\right)} \qquad\Longleftrightarrow\qquad \Delta_rG^\circ + RT\ln K^\circ = 0. \]
Elle est sans dimension et ne dépend que de $T$ (et de l'écriture de l'équation).
\end{definition}

\begin{resultat}{Loi d'action des masses (Guldberg et Waage) et critère en $Q_r$}
En combinant avec l'expression de $\Delta_rG$ :
\[ \boxed{\Delta_rG = RT\ln\frac{Q_r}{K^\circ}} \]
\begin{itemize}
  \item à l'équilibre, $\Delta_rG = 0$ donc $Q_{r,\text{éq}} = K^\circ(T)$ : c'est la \textbf{loi d'action des masses} ;
  \item hors équilibre, $Q_r < K^\circ \Leftrightarrow \Delta_rG < 0$ : sens direct ; $Q_r > K^\circ \Leftrightarrow \Delta_rG > 0$ : sens indirect.
\end{itemize}
\end{resultat}

\begin{figure}[H]
\centering
\begin{tikzpicture}[>=Stealth]
  \draw[->,line width=1pt] (0,0) -- (11,0) node[right] {$Q_r$};
  \draw (0,0.12) -- (0,-0.12) node[below] {$0$};
  \draw[line width=1.2pt,c3] (5.5,0.25) -- (5.5,-0.25) node[below,black] {$K^\circ$};
  \draw[->,line width=1.4pt,c4] (2.6,0.55) -- (4.6,0.55);
  \node[c4,font=\small,align=center] at (2.75,1.25) {$Q_r<K^\circ$ : $\Delta_rG<0$\\évolution dans le sens direct};
  \draw[->,line width=1.4pt,c2] (8.4,0.55) -- (6.4,0.55);
  \node[c2,font=\small,align=center] at (8.25,1.25) {$Q_r>K^\circ$ : $\Delta_rG>0$\\évolution dans le sens indirect};
  \node[c3,font=\small] at (5.5,-1.05) {équilibre : $Q_r = K^\circ$, $\Delta_rG = 0$};
\end{tikzpicture}
\caption{Le quotient de réaction évolue toujours vers la constante d'équilibre.}
\label{fig:axeQK}
\end{figure}

\begin{resultat}{Ce que dit le signe de $\Delta_rG^\circ$}
$K^\circ > 1 \Leftrightarrow \Delta_rG^\circ < 0$ : la réaction est dite \emph{thermodynamiquement favorisée} dans le sens direct (à partir de constituants dans l'état standard, $Q_r = 1 < K^\circ$). En pratique, si $K^\circ \gtrsim 10^4$, la réaction est souvent quasi totale ; si $K^\circ \lesssim 10^{-4}$, elle est quasi nulle. Ce ne sont que des repères : l'état final dépend aussi des quantités initiales.
\end{resultat}

\begin{resultat}{Les visages de $\Delta_rG^\circ$ -- à connaître}
\[ \Delta_rG^\circ = \sum_i\nu_i\mu_i^\circ = \sum_i\nu_i\,\Delta_fG^\circ_i = \Delta_rH^\circ - T\,\Delta_rS^\circ = -RT\ln K^\circ \]
et, pour une réaction rédox (chapitre~\ref{ch:redox}), $\Delta_rG^\circ = -\gamma F\,\Delta E^\circ$.
\end{resultat}

\begin{exemple}[ammoniac à $\SI{298}{K}$]
$\Delta_rG^\circ = \SI{-32,7}{kJ.mol^{-1}}$ (chapitre~\ref{ch:grandeurs}), donc
\[ K^\circ(\SI{298}{K}) = \exp\left(\frac{32{,}7\times10^3}{8{,}314\times298}\right) = \num{5,4e5}. \]
La synthèse est très favorisée à température ambiante\dots\ mais infiniment lente sans catalyseur ni chauffage (chapitre~\ref{ch:vanthoff}).
\end{exemple}

\section{Représentation graphique : la courbe $G(\xi)$}

À $T$ et $P$ fixées, $G$ ne dépend plus que de $\xi$. On peut la calculer avec $G = \sum_in_i\mu_i$ et le tableau d'avancement. La pente de la courbe est $\left(\partial G/\partial\xi\right)_{T,P} = \Delta_rG$.

\begin{exemple}[isomérisation en phase gazeuse]
\ce{A(g) = B(g)} sous $P = P^\circ$, avec $n_A = 1-\xi$, $n_B = \xi$ (en mol), donc $a_A = 1-\xi$ et $a_B = \xi$. En prenant $\mu_A^\circ = 0$ comme référence, $\mu_B^\circ = \Delta_rG^\circ$ :
\[ G(\xi) = \underbrace{\xi\,\Delta_rG^\circ}_{\text{réactifs et produits séparés}} + \underbrace{RT\big[(1-\xi)\ln(1-\xi) + \xi\ln\xi\big]}_{\text{terme de mélange}<0} \]
Avec $\Delta_rG^\circ = \SI{-2,0}{kJ.mol^{-1}}$ à $\SI{298}{K}$ : $K^\circ = 2{,}24$ et $\xi_{\text{éq}} = K^\circ/(1+K^\circ) = \SI{0,69}{mol}$.
\end{exemple}

\begin{figure}[H]
\centering
\begin{tikzpicture}
\begin{axis}[cours,width=0.76\linewidth,xlabel={$\xi$ (mol)},ylabel={$G - G(0)$ (kJ)},xmin=0,xmax=1.02,ymin=-3.3,ymax=0.4,
  legend style={at={(0.97,0.97)},anchor=north east}]
  \addplot[c1] table[x=xi,y=G]{data/gxi.dat}; \addlegendentry{$G(\xi)$}
  \addplot[black!50,dashed,line width=0.9pt] table[x=xi,y=Glin]{data/gxi.dat}; \addlegendentry{$\xi\,\Delta_rG^\circ$ (sans mélange)}
  \addplot[c4,line width=1.3pt] coordinates {(0.15,-1.499) (0.45,-2.729)};
  \addplot[c2,line width=1.3pt] coordinates {(0.80,-2.951) (1.00,-2.261)};
  \addplot[only marks,mark=*,mark size=1.8pt,c4] coordinates {(0.3,-2.114)};
  \addplot[only marks,mark=*,mark size=1.8pt,c2] coordinates {(0.9,-2.606)};
  \addplot[c3,dashed,line width=0.8pt] coordinates {(0.691,-3.3) (0.691,-2.915)};
  \addplot[only marks,mark=*,mark size=2.2pt,c3] coordinates {(0.691,-2.915)};
  \node[c4,font=\small,anchor=west] at (axis cs:0.02,-2.85) {$\Delta_rG<0$ : $\xi$ augmente};
  \node[c2,font=\small,anchor=east] at (axis cs:0.93,-2.18) {$\Delta_rG>0$ : $\xi$ diminue};
  \node[c3,font=\small,anchor=north east] at (axis cs:0.68,-3.0) {$\xi_{\text{éq}}$ : $\Delta_rG = 0$};
\end{axis}
\end{tikzpicture}
\caption{Courbe $G(\xi)$ : le système « descend » vers le minimum. La pente de la tangente est $\Delta_rG$.}
\label{fig:gxi}
\end{figure}

\noindent\textbf{Pourquoi le minimum est-il à l'intérieur ?} Le terme de mélange a une pente infinie en $\xi\to0$ et $\xi\to1$ ($RT\ln\xi\to-\infty$) : dès qu'un constituant d'une phase mélange (gaz, soluté) a une quantité nulle, $\Delta_rG$ est infini et pousse le système à le former. Un équilibre entre gaz ou solutés n'est donc jamais rigoureusement total. Il en va autrement si les seuls constituants qui disparaissent sont des solides ou liquides \emph{purs} : leur activité reste égale à 1 jusqu'à leur disparition, et la réaction peut s'arrêter par épuisement avant d'atteindre l'équilibre.

\section{État final : équilibre ou transformation totale}

\begin{definition}[équilibre chimique, transformation totale, rupture d'équilibre]
\begin{itemize}
  \item L'état final est un \textbf{équilibre chimique} si tous les constituants de l'équation sont présents : alors $\Delta_rG = 0$ et $Q_r = K^\circ$.
  \item La transformation est \textbf{totale} si un réactif est entièrement consommé : $\xi_F = \xi_{\max}$. Le système s'arrête alors \emph{hors équilibre}, avec $\Delta_rG < 0$ ($Q_r < K^\circ$), faute de réactif.
  \item On parle de \textbf{rupture d'équilibre} quand l'espèce épuisée est une phase condensée pure (solide, liquide), ou un précipité qui se dissout complètement.
\end{itemize}
\end{definition}

\begin{methode}[déterminer l'état final]
\begin{enumerate}
  \item Tableau d'avancement ; calculer $\xi_{\max}$ (et $\xi_{\min}$ si le sens n'est pas connu).
  \item Calculer $Q_{r,0}$ et le comparer à $K^\circ$ pour connaître le sens d'évolution.
  \item \emph{Supposer} l'équilibre atteint : résoudre $Q_r(\xi_{\text{éq}}) = K^\circ$.
  \item \emph{Vérifier} : si $\xi_{\text{éq}}$ est accessible ($\xi_{\text{éq}}\leq\xi_{\max}$), l'hypothèse est validée. Sinon, la transformation est totale : $\xi_F = \xi_{\max}$, et l'on vérifie $Q_r(\xi_{\max}) < K^\circ$.
  \item Si $K^\circ$ est très grande, on peut supposer d'emblée la réaction quasi totale, puis calculer la très faible quantité restante du réactif limitant à l'aide de $K^\circ$.
\end{enumerate}
\end{methode}

\begin{exemple}[décomposition du calcaire en enceinte fermée]
\ce{CaCO3(s) = CaO(s) + CO2(g)}, $K^\circ = P_{\ce{CO2}}/P^\circ$ (les solides ont une activité 1). Données : $\Delta_rH^\circ = \SI{179,2}{kJ.mol^{-1}}$, $\Delta_rS^\circ = \SI{160,2}{J.K^{-1}.mol^{-1}}$.

\noindent À $T = \SI{1100}{K}$ : $\Delta_rG^\circ = 179{,}2 - 1100\times0{,}1602 = \SI{3,0}{kJ.mol^{-1}}$ et $K^\circ = \mathrm{e}^{-3{,}0\times10^3/(8{,}314\times1100)} = 0{,}722$.

\noindent On introduit $n_0$ de \ce{CaCO3} dans une enceinte vide de $V = \SI{10}{L}$. Si l'équilibre est atteint, $P_{\ce{CO2}} = \SI{0,722}{bar}$, ce qui demande
\[ \xi_{\text{éq}} = \frac{P_{\ce{CO2}}V}{RT} = \frac{0{,}722\times10^5\times10^{-2}}{8{,}314\times1100} = \SI{78,9}{mmol}. \]
\begin{itemize}
  \item $n_0 = \SI{0,100}{mol} > \xi_{\text{éq}}$ : équilibre, il reste $\SI{21}{mmol}$ de \ce{CaCO3}, $P = \SI{0,722}{bar}$.
  \item $n_0 = \SI{0,050}{mol} < \xi_{\text{éq}}$ : \textbf{rupture d'équilibre}. Tout le calcaire se décompose ; $P = \dfrac{n_0RT}{V} = \SI{0,457}{bar} < K^\circ P^\circ$, et $\Delta_rG = RT\ln(Q_r/K^\circ) = \SI{-4,2}{kJ.mol^{-1}} < 0$ : la réaction « voudrait » continuer mais il n'y a plus de réactif.
\end{itemize}
\end{exemple}

\begin{figure}[H]
\centering
\begin{tikzpicture}
\begin{axis}[cours,width=0.72\linewidth,height=0.4\linewidth,xlabel={$n_0(\ce{CaCO3})$ (mmol)},ylabel={$P_{\ce{CO2}}$ final (bar)},
  xmin=0,xmax=150,ymin=0,ymax=0.9]
  \addplot[c1] table[x=n0_mmol,y=P_bar]{data/caco3.dat};
  \addplot[c3,dashed,line width=0.8pt] coordinates {(78.9,0) (78.9,0.722)};
  \node[c3,font=\small,anchor=north west] at (axis cs:80,0.18) {$\xi_{\text{éq}} = \SI{78,9}{mmol}$};
  \node[font=\small,align=center,c2] at (axis cs:36,0.68) {transformation\\totale};
  \node[font=\small,align=center,c4] at (axis cs:117,0.55) {équilibre\\$P = K^\circ P^\circ$};
\end{axis}
\end{tikzpicture}
\caption{Pression finale en \ce{CO2} à $\SI{1100}{K}$ dans $\SI{10}{L}$ en fonction de la quantité initiale de calcaire.}
\label{fig:caco3}
\end{figure}
